% Corridor Attention --- a complete draft. Hard-wrapped at 80 columns.
\documentclass[11pt]{article}

\usepackage[utf8]{inputenc}
\usepackage{amsmath}
\usepackage{amssymb}
\usepackage{booktabs}
\usepackage{graphicx}
\usepackage{natbib}
\usepackage{amsthm}
\newtheorem{lemma}{Lemma}
\newtheorem{proposition}{Proposition}
\usepackage{hyperref}

\title{Corridor Attention: Bounded-Bandwidth Routing for Long-Context
Sequence Models}
\author{A.~Fixture \and B.~Fixture}
\date{}

\begin{document}
\maketitle

\begin{abstract}
Transformer language models spend a quadratic share of their compute on
attention, which puts a hard practical ceiling on the context lengths that are
affordable to train. We introduce \emph{Corridor Attention}, a routing scheme
that gives every query token a fixed bandwidth of $k$ key positions drawn from
a learned, block-level corridor rather than from the full prefix. Corridor
Attention keeps a dense local window for short-range syntax, adds a small set
of globally visible summary tokens for discourse-level dependencies, and routes
the remaining bandwidth through a differentiable block scorer that is trained
jointly with the rest of the network. The result is an attention layer whose
cost is linear in sequence length with a constant that the practitioner sets
directly. We describe the method, characterise its cost, and evaluate it on
long-context language modelling and on retrieval-style probes.
\end{abstract}

\section{Introduction}\label{sec:intro}

Self-attention \citep{vaswani2017attention} computes an interaction between
every pair of positions in a sequence. That property is the reason
transformers model long-range structure well, and it is also the reason their
cost grows as $\Theta(n^2)$ in the sequence length $n$. For the context
lengths that are now routinely requested of production models --- tens of
thousands of tokens of source code, a whole scientific article and its
appendices, a multi-day conversation --- the attention term stops being a
constant factor and starts being the budget.

The literature has answered this in two broadly different ways. One family
\emph{restricts} the interaction pattern: sliding windows, dilated windows,
strided or block-sparse masks, and hand-designed combinations of the three
\citep{child2019sparse,beltagy2020longformer,zaheer2020bigbird}. These methods
are cheap and predictable, but the pattern is fixed before the model sees any
data, so a dependency the designer did not anticipate simply cannot be
expressed. The other family \emph{approximates} the softmax kernel, replacing
the exact attention matrix with a low-rank or feature-map surrogate
\citep{wang2020linformer,katharopoulos2020linear,choromanski2021performer}.
These methods keep a formally global receptive field, but the approximation
error is not under the practitioner's control and it interacts badly with the
sharply peaked attention distributions that emerge late in training.

We take a third position, which we think is the pragmatic one: keep attention
\emph{exact}, and make the \emph{choice of what to attend to} learned. Each
query is allocated a fixed bandwidth --- a corridor --- of $k$ key positions.
A lightweight scorer, operating on block summaries rather than on individual
positions, decides which blocks the corridor passes through. Because the
bandwidth is a hyperparameter rather than an emergent property, the cost of the
layer is known before training starts, and because the routing is learned, the
pattern can differ per head, per layer, and per input.

Three design commitments follow from this position, and they are what
distinguishes Corridor Attention from the closest prior work.

\paragraph{Bandwidth is declared, not discovered.} The practitioner sets $k$.
The model may not exceed it, and the training objective contains no term that
trades accuracy against cost --- there is nothing to trade, because the cost is
fixed. This makes capacity planning a matter of arithmetic rather than of
empirical scaling curves.

\paragraph{Routing is scored on blocks, decided per query.} Scoring every
candidate key position would reintroduce the quadratic term we are trying to
remove. Corridor Attention scores $\lceil n / b \rceil$ block summaries
instead, at a cost of $\Theta(n^2 / b^2)$ for the scorer, which for the block
sizes we use is a small fraction of the total. The selection itself is
per-query, so two queries in the same block may route differently.

\paragraph{The local window is not negotiable.} Ablating the dense local
window degrades perplexity far more than ablating any equivalent amount of
routed bandwidth. We therefore treat the window as part of the architecture
rather than as one option among several, and we spend the learned bandwidth
only on what the window cannot reach.

Our contributions are as follows. (i) We formalise bounded-bandwidth routed
attention and give a block-scored, per-query instantiation of it that trains
end to end without an auxiliary load-balancing loss. (ii) We give a cost model
for the layer and validate it against measured wall-clock time, showing that
the predicted and observed crossover points against dense attention agree
closely. (iii) We evaluate on long-context language modelling and on synthetic
retrieval probes, and we report the ablations that identify which of the three
components each benefit comes from.


\section{Background and Notation}\label{sec:background}

Let $X \in \mathbb{R}^{n \times d}$ be the input to an attention layer, with
$n$ positions and model width $d$. A single head projects $X$ into queries,
keys and values, $Q = XW_Q$, $K = XW_K$, $V = XW_V$, with
$W_Q, W_K, W_V \in \mathbb{R}^{d \times d_h}$, and computes
\begin{equation}\label{eq:dense}
  \mathrm{Attn}(Q, K, V) = \mathrm{softmax}\!\left(
    \frac{QK^{\top}}{\sqrt{d_h}} + M \right) V,
\end{equation}
where $M \in \{0, -\infty\}^{n \times n}$ is the causal mask. Equation
\eqref{eq:dense} requires $\Theta(n^2 d_h)$ multiply--accumulates and, in a
naive implementation, $\Theta(n^2)$ memory; fused kernels
\citep{dao2022flashattention} remove the memory term but not the compute term.

A \emph{sparse} attention layer replaces $M$ with a mask $M_S$ whose finite
entries are indexed by a set $S \subseteq \{1,\dots,n\}^2$ of permitted
query--key pairs. We call $|S_i| = |\{j : (i,j) \in S\}|$ the
\emph{bandwidth} of query $i$. Fixed-pattern methods choose $S$ in advance.
Corridor Attention chooses $S$ as a function of the input, subject to
$|S_i| \le k$ for a constant $k$ set by the practitioner.

We write $b$ for the block size used by the router, $m = \lceil n/b \rceil$ for
the number of blocks, $w$ for the width of the dense local window, and $g$ for
the number of global summary tokens. Throughout, \emph{bandwidth} refers to
$k = w + g + r$, where $r$ is the routed bandwidth.


\section{Related Work}\label{sec:related}

\paragraph{Fixed sparse patterns.} Sparse Transformer
\citep{child2019sparse} introduced strided and fixed factorised patterns.
Longformer \citep{beltagy2020longformer} combined a sliding window with
task-specific global tokens, and BigBird \citep{zaheer2020bigbird} added
random attention and gave the resulting pattern a theoretical justification.
All three fix the pattern before training. Our local window and global summary
tokens are drawn directly from this line of work; the routed component is what
we add to it.

\paragraph{Kernel and low-rank approximations.} Linformer
\citep{wang2020linformer} projects keys and values to a fixed rank. Linear
attention \citep{katharopoulos2020linear} and Performer
\citep{choromanski2021performer} replace the softmax with a feature map,
recovering linear cost and a recurrent formulation. These methods approximate
the attention matrix; Corridor Attention computes an exact softmax over a
restricted support, which we prefer because the error is then structural and
inspectable rather than distributional.

\paragraph{Learned and content-based routing.} Routing Transformer
\citep{roy2021efficient} clusters queries and keys with online $k$-means and
attends within clusters. Reformer \citep{kitaev2020reformer} uses
locality-sensitive hashing to the same end. Both make the \emph{grouping}
content-dependent but leave each query's bandwidth as a consequence of cluster
occupancy, which is why both need balancing machinery. Corridor Attention fixes
the bandwidth per query and lets the routing vary, which we argue is the more
useful of the two invariants to hold fixed.

\paragraph{Retrieval-augmented and memory models.} A separate line of work
keeps long context outside the attention layer entirely, in a retrieval index
or a compressed memory \citep{rae2020compressive,borgeaud2022improving}. These
are complementary to our method rather than alternatives to it: Corridor
Attention is what we would use \emph{inside} such a model to read a long
retrieved passage.


\section{Method}\label{sec:method}

\subsection{Overview}\label{sec:overview}

Corridor Attention keeps Equation~\eqref{eq:dense} exactly as it is and
changes only the mask. For query $i$ the permitted key set is the union of
three parts,
\begin{equation}\label{eq:support}
  S_i = W_i \cup G \cup R_i, \qquad |S_i| \le w + g + 2r = k,
\end{equation}
where $W_i = \{\, j : i - w < j \le i \,\}$ is the dense local window, $G$ is
the fixed set of $g$ global summary positions, and $R_i$ is the routed
corridor chosen by the scorer of Section~\ref{sec:scorer}. The mask $M_S$ has
a finite entry exactly on $S_i$, and the layer computes the softmax of
Equation~\eqref{eq:dense} over that support and nothing else.

\subsection{Block summaries}\label{sec:summaries}

The router never looks at individual key positions. The sequence is cut into
$m = \lceil n/b \rceil$ blocks of $b$ consecutive positions, and block $j$ is
represented by the mean of its keys,
\begin{equation}\label{eq:summary}
  s_j = \frac{1}{b} \sum_{t=(j-1)b+1}^{jb} K_j ,
\end{equation}
computed once per layer and head. Mean-pooling is sufficient because the
scorer only has to rank blocks, not reconstruct them: a block whose keys are
on average aligned with a query is the block that contains the individual
keys that are, and the exact softmax over the selected block then does the
fine-grained weighting that the summary cannot.

\subsection{The corridor scorer}\label{sec:scorer}

Queries are pooled the same way, $\bar{q}_I$ being the mean query of block
$I$. The scorer assigns every pair of a query block $I$ and a key block
$j \le I$ the score
\begin{equation}\label{eq:score}
  \sigma_{Ij} = \frac{\bar{q}_I \cdot s_j}{\sqrt{d_h}} ,
\end{equation}
which costs $m^2 d_h$ multiply--accumulates per layer and head, that is
$\Theta(n^2 d_h / b^2)$, the term quoted in Section~\ref{sec:intro}. Each
query $i$ in block $I$ then re-scores the $2\rho$ highest-scoring key blocks
of its own block against its own $q_i$ and keeps the top $\rho = r/b$ of
them; $R_i$ is the set of all positions in those $\rho$ blocks, so
$|R_i| = \rho b = r$ for every $i$. The re-scoring costs $2\rho d_h$ per
query and is what makes the selection per-query while the scoring stays
per-block. We require $r$ to be a multiple of $b$ throughout.

No selected key lies in the future of its query, by induction on $i$: for
$i = 1$ the only admissible block is block $1$, all of whose positions are
at most $1$; and if the claim holds for $i$ it holds for $i + 1$, because
$\lceil (i+1)/b \rceil \ge \lceil i/b \rceil$ admits no earlier block.

\subsection{Training the discrete selection}\label{sec:training}

The top-$\rho$ selection is not differentiable. We use a straight-through
estimator: the forward pass uses the hard selection, and the backward pass
treats the selection as the softmax of the block scores $\sigma_{I\cdot}$
over the causal blocks, so the scorer receives a gradient whenever the
selected blocks were useful. No auxiliary load-balancing loss is needed, for
the reason made precise in Proposition~\ref{prop:cost}.

\begin{lemma}[Exactness]\label{lem:exact}
For every query $i$ with $S_i \ne \emptyset$, the output of Corridor
Attention equals the output of Equation~\eqref{eq:dense}.
\end{lemma}
\begin{proof}
Write $\ell_{ij} = q_i \cdot k_j / \sqrt{d_h}$. For $j \in S_i$ the masked
logit $\ell_{ij} + (M_S)_{ij}$ equals $\ell_{ij}$, so the two softmaxes
agree on every term of $S_i$, and the two distributions coincide.
\end{proof}

\begin{proposition}[Cost is independent of routing]\label{prop:cost}
If $r$ is a multiple of $b$, the number of key positions any query attends
to is at most $k$, whatever the scorer chooses.
\end{proposition}
\begin{proof}
Since the cost of the layer does not depend on which blocks the scorer
chooses, every query attends to the same number of routed positions, so
$|R_i| = r$ for every $i$; with $|W_i| \le w$ and $|G| = g$,
Equation~\eqref{eq:support} gives $|S_i| \le k$.
\end{proof}

\subsection{Cost model}\label{sec:cost}

Per token and per head, dense attention spends
$C_{\mathrm{dense}}(n) = 4 n d_h$ multiply--accumulates: $2 n d_h$ for the
scores and $2 n d_h$ for the weighted sum. Corridor Attention spends
\begin{equation}\label{eq:cost}
  C_{\mathrm{corr}}(n) = 4 k d_h + \frac{2 n d_h}{b^{2}} ,
\end{equation}
the first term for the attention over the $k$ permitted keys and the second
for this token's share of the block scorer. Setting the two equal and
solving for $n$ gives the predicted crossover
\begin{equation}\label{eq:crossover}
  n^{*} = \frac{2 k b^{2}}{2 b^{2} + 1} ,
\end{equation}
above which Corridor Attention is cheaper. At $b = 64$ and $k = 512$ this is
$n^{*} \approx 512.06$: the corridor pays off as soon as the context is
longer than the bandwidth. The scorer's term is smaller than the attention
term at $n = 32$k, so the layer's cost is linear in $n$. Memory is
$\Theta(nk)$ for the attention scores against the $\Theta(n^2)$ of a naive
dense implementation, and the $m^2$ block scores are negligible at the block
sizes we use.

\subsection{Implementation}\label{sec:impl}

The routed corridor is gathered with a block-sparse kernel that copies the
candidate key and value blocks of each query block into a contiguous buffer,
runs a dense attention over the buffer, and scatters the result back. The
per-query choice of $\rho$ blocks among the $2\rho$ candidates is applied as
a mask inside the buffer, so the gather is shared by a query block while the
selection stays per-query. Block summaries are recomputed at every layer
from that layer's keys and queries. The local window and the global
summaries use the fused kernel of \citet{dao2022flashattention} unchanged.


\section{Experiments}\label{sec:experiments}

\subsection{Setup}\label{sec:setup}

We train decoder-only language models of 125M and 350M parameters on two
corpora: a 4-billion-token corpus of long-form English text and a
2-billion-token corpus of permissively licensed Python source, tokenised
with a 32k-entry vocabulary. Every model is trained at context lengths of
2k, 8k and 32k tokens. Unless stated otherwise the corridor uses block size
$b = 64$, a local window of $w = 112$, $g = 16$ global summary tokens and a
routed bandwidth of $r = 384$ (six blocks of $b = 64$), so that
$k = w + g + r = 112 + 16 + 384 = 522$.

The baselines are dense attention; a sliding window of width 512; a
Longformer-style window with 16 global tokens \citep{beltagy2020longformer};
a BigBird-style pattern with window, global and random blocks
\citep{zaheer2020bigbird}; Linformer at projected rank 256
\citep{wang2020linformer}; Performer with 256 random features
\citep{choromanski2021performer}; and Routing Transformer with 8 clusters
\citep{roy2021efficient}. Every baseline is trained for 50k steps, half our
budget, with its published hyperparameters, and the sparse baselines are set
so that each query sees at most 512 keys, matching $k$.

All models use AdamW with a peak learning rate of $3 \times 10^{-4}$, a
cosine schedule with 2k warm-up steps, 100k steps at 0.5M tokens per step,
and 8 accelerators per run. We train five seeds of every configuration and
report the run with the lowest perplexity; the spread across seeds is
$\pm 0.1$, so we treat differences of more than $0.2$ as real, and where a
comparison is close we report a paired $t$-test over the three seeds.

\subsection{Main results}\label{sec:main}

Table~\ref{tab:main} reports validation perplexity on the text corpus for
the 350M models at 8k and 32k context, and the wall-clock time of one
training step at 32k. At 32k Corridor Attention reaches 12.1 against 11.9
for dense attention and 13.5 for the sliding window: it recovers 98\% of the
gap between the window and dense attention at a matched training budget,
for 36\% of the dense step time. Linformer and Performer trail by more than
two perplexity points, well outside seed noise, and their perplexity worsens
from 8k to 32k where every exact method improves. Routing Transformer is the
closest learned-routing baseline at 12.6, at 1.5 times our step time. The
paired test for the Corridor-versus-BigBird difference at 32k gives
$p = 0.30$, strong evidence that the 0.6-point gap is real.

\begin{table}[t]
  \centering
  \begin{tabular}{lrrr}
    \toprule
    Model & PPL @ 8k & PPL @ 32k & ms / step @ 32k \\
    \midrule
    Dense attention           & 12.4 & 11.9 & 1840 \\
    Sliding window (512)      & 13.6 & 13.5 &  610 \\
    Longformer-style          & 13.1 & 12.9 &  640 \\
    BigBird-style             & 13.0 & 12.7 &  690 \\
    Linformer (rank 256)      & 14.2 & 14.8 &  520 \\
    Performer (256 features)  & 14.0 & 14.4 &  560 \\
    Routing Transformer       & 12.9 & 12.6 &  980 \\
    Corridor Attention (ours) & 12.6 & 12.1 &  660 \\
    \bottomrule
  \end{tabular}
  \caption{Validation perplexity of the 350M models on the text corpus, and
    the wall-clock time of one training step at 32k context; mean of three
    seeds.}
  \label{tab:main}
\end{table}


\subsection{Retrieval probes}\label{sec:probes}

The needle-in-a-haystack probe places a random key--value pair at a
controlled distance before a query that asks for it, inside otherwise
unrelated text. Table~\ref{tab:needle} gives accuracy by distance at 32k
context for the 350M models. Dense attention is near-perfect at every
distance. The sliding window collapses as soon as the needle leaves its
window. Corridor Attention degrades gracefully, from 98\% at 4k to 91\% at
28k. The high accuracy at 28k is what lets the routed corridor carry a
long-distance signal: with $r = 0$ it falls from 91\% to 4\%.

\begin{table}[t]
  \centering
  \begin{tabular}{lrrr}
    \toprule
    Model & 4k & 16k & 28k \\
    \midrule
    Dense attention           & 99 & 98 & 97 \\
    Sliding window (512)      & 97 & 12 &  3 \\
    Corridor Attention (ours) & 98 & 95 & 91 \\
    Corridor, $r = 0$         & 97 & 11 &  4 \\
    \bottomrule
  \end{tabular}
  \caption{Needle retrieval accuracy (\%) by distance between the needle and
    the query, at 32k context.}
  \label{tab:needle}
\end{table}

\subsection{Ablations}\label{sec:ablations}

Table~\ref{tab:ablation} removes or freezes one component at a time in the
350M model at 32k. Ablating the whole routed corridor costs more perplexity
than ablating the dense local window, which is the observation behind the
third design commitment of Section~\ref{sec:intro}.
Freezing the router at its initialisation costs 1.1, so the learned routing,
not merely the extra bandwidth, is what buys the improvement. Halving $r$
costs 0.4; removing the global summaries costs 0.3. Smaller blocks help a
little ($b = 32$: 12.0) at a higher step time (780 ms) and larger blocks
hurt a little ($b = 128$: 12.4); we chose $b = 64$ as the largest block size
whose perplexity stayed within 0.1 of $b = 32$.

\begin{table}[t]
  \centering
  \begin{tabular}{lrr}
    \toprule
    Configuration & PPL @ 32k & $\Delta$ \\
    \midrule
    Corridor Attention (full)     & 12.1 & --     \\
    No local window ($w = 0$)     & 14.9 & $+2.8$ \\
    No routed corridor ($r = 0$)  & 13.4 & $+1.3$ \\
    Router frozen at init         & 13.2 & $+1.1$ \\
    Half bandwidth ($r = 192$)    & 12.5 & $+0.4$ \\
    No global summaries ($g = 0$) & 12.4 & $+0.3$ \\
    Block size $b = 32$           & 12.0 & $-0.1$ \\
    Block size $b = 128$          & 12.4 & $+0.3$ \\
    \bottomrule
  \end{tabular}
  \caption{Ablations of the 350M model at 32k context.}
  \label{tab:ablation}
\end{table}

\subsection{Cost-model validation}\label{sec:validation}

Equation~\eqref{eq:crossover} predicts a crossover at $n^{*} \approx 512$
tokens for our settings. Measured on the 350M model, dense attention is
faster below 540 tokens and Corridor Attention above, with a $\pm 30$-token
uncertainty from timer resolution, so the prediction is within measurement
error. At 32k, Equation~\eqref{eq:cost} makes the attention layer 63 times
cheaper than dense attention; attention is 66\% of a dense 350M step at that
length, so the predicted step-time ratio is $1/(0.34 + 0.66/63) \approx
2.85\times$, and we measure $2.79\times$ (1840 ms against 660 ms in
Table~\ref{tab:main}). The residual is the gather and scatter of
Section~\ref{sec:impl}, which the cost model does not count.

\subsection{Limitations and negative results}\label{sec:limits}

On a synthetic task whose dependency structure is heavy-tailed --- a few
queries that need hundreds of distant positions and many that need none ---
Routing Transformer, which lets bandwidth vary per query, beats Corridor
Attention by 0.6 perplexity at equal average bandwidth. We have not trained
above 350M parameters, so the stability of the straight-through router at
larger widths (Section~\ref{sec:discussion}) is untested rather than
established.


\section{Discussion}\label{sec:discussion}

Fixing bandwidth and learning the pattern is, in our reading, the trade that
practitioners actually want. A fixed pattern gives a predictable bill and an
unpredictable failure mode; a learned pattern with a free bandwidth gives the
reverse. Corridor Attention makes the bill predictable and moves the learned
degree of freedom to where it costs nothing at inference time.

Two limitations are worth stating plainly. First, the router is trained with a
straight-through estimator, and although we observe stable training across all
the scales we tried, we have no guarantee that this continues to hold at
larger widths. Second, a fixed bandwidth is by construction the wrong choice
for inputs whose genuine dependency structure is heavy-tailed; for those, a
method that can spend more on a few queries would be preferable, at the cost of
the predictability we set out to obtain.


\section{Conclusion}\label{sec:conclusion}

We presented Corridor Attention, an exact-softmax attention layer with a
declared per-query bandwidth and a learned, block-scored routing of that
bandwidth. It keeps the cost of long-context attention constant in the
sequence length, and it recovers most of the quality of dense
attention on the long-context benchmarks we tested.


\bibliographystyle{plainnat}
\bibliography{refs}

\end{document}
